\documentclass{beamer}
\usefonttheme{professionalfonts}
%\usetheme{Rochester}
%\usetheme{AnnArbor}
\usetheme{Szeged}
\useoutertheme{infolines}
\usecolortheme{dolphin}
\usepackage{graphicx}
% \usepackage[english]{babel}
\usepackage{hyperref}
\usepackage{multicol}
\usepackage{color}
\usepackage{xcolor}
%\usepackage{smartdiagram}


%\definecolor{carmine}{rgb}{0.59, 0.0, 0.09}
\usepackage{amsmath} 

\usepackage{graphicx}
\usepackage{tikz}
\usetikzlibrary{shapes.callouts}
\usepackage{multirow}

\usepackage[utf8]{inputenc}

\usepackage{enumerate}
\usepackage[shortlabels]{enumitem}
\usepackage{amsmath,amsthm,amssymb,amsfonts,mathrsfs,latexsym,stmaryrd}
\usepackage{tcolorbox}

\usepackage{color} 
\usepackage{amssymb, amsmath, amsbsy}
\usepackage{adjustbox,lipsum}
\usepackage{placeins}

\usepackage{textpos}
\usepackage{wrapfig}
\usepackage[spanish, es-tabla]{babel}

\usepackage{caption}
\setlength{\parindent}{0pt}
\usepackage[hyperfootnotes=false]{hyperref}

\usepackage{url}
\usepackage{float}
\usepackage{tikz}
\usetikzlibrary{patterns}


\title[ICS-162]{Econometría ICS-294\\ \small{Clase 20: Raíces unitarias}}

\author{Jocelyn Tapia Stefanoni}
\institute[USM]
{Universidad Técnica Federico Santa María\\ % Your institution for the title page
	\medskip
	\textit{} % Your email address
	
}
\logo{

\includegraphics[width=3cm]{figuras/logoUSM.png}
}
\date{}
\newcommand\Fontvi{\fontsize{9}{7.2}\selectfont}
\begin{document}
	
\begin{frame}
	\titlepage % Print the title page as the first slide


\end{frame}
%%%%%%%%%%%%%%%%%%%%%%%%%%%%%%%%%%%%%%%%%%%%%%%%%%%%%%%%%%%%%%%%%%%%%%%%%%%
\logo{}

\begin{frame}{Introducción}
    \begin{itemize}
        \item La clase pasada vimos datos de series de tiempo.
        \item Estos tienen la particularidad de estar ordenados y pueden presentar tendencias temporales.
        \item Esto es un problema a la hora de interpretar las estimaciones ya que podemos estar frente a \textbf{correlaciones espurias}:
        \begin{itemize}
            \newline -- Correlaciones muy fuertes entre dos variables que no se deben más que a tendencias temporales en cada una de ellas.
            \newline -- Si la tendencia es determinística se puede solucionar incluyendo una tendencia temporal como regresor en el MRL.
        \end{itemize}
        \item También vimos que hay series con tendencias estocásticas \textbf{altamente persistentes }–caminatas aleatorias– para las que precisamos diferenciar para poder estimar.
        \item Hoy vamos a ver como testear la “alta persistencia” con test de raíces unitarias de Dickey-Fuller.
    \end{itemize}
\end{frame}
\begin{frame}{Raíces unitarias, tendencia estocástica y regresión espuria}
    \begin{itemize}
        \item Queremos estimar el modelo:
        \[
        Y_t = \beta_0 + \beta_1 X_t + u_t
        \]
        \item Supongamos que el proceso de cada variable es el siguiente:
        \[
        Y_t = \alpha_1 + \rho_1 Y_{t-1} + \epsilon_t
        \]
        \[
        X_t = \alpha_2 + \rho_2 X_{t-1} + v_t
        \]
        \item Si $-1 < \rho < 1$, entonces los procesos son AR(1) estacionarios y se puede hacer la regresión sin problema. También se los llama integrados de orden 0.
        \item Si $\rho = 1$ la serie es una caminata aleatoria. Diremos que tiene raíz unitaria o que es integrada de orden 1.
        \item Esto es problemático y hace que se encuentren relaciones espurias entre $x$ e $y$: $\beta_1$ significativo cuando en realidad no hay relación.
    \end{itemize}
\end{frame}
\begin{frame}{Solución al problema de regresión espuria por raíz unitaria: diferenciar las variables}
    \begin{itemize}
        \item Si $\rho = 1$, en lugar de estimar el modelo original, hacemos lo siguiente:
        \begin{enumerate}
            \item Diferenciamos las variables:
            \[
            \Delta Y_t = Y_t - Y_{t-1}, \quad \Delta X_t = X_t - X_{t-1}
            \]
            \item Hacemos la regresión en diferencias:
            \[
            \Delta Y_t = \beta_0 + \beta_1 \Delta X_t + u_t
            \]
        \end{enumerate}
        \item Intuición:
        \[
        Y_t = \alpha_1 + Y_{t-1} + \epsilon_t \implies \text{altamente persistente}
        \]
        \[
        \Delta Y_t = \alpha_1 + \epsilon_t \implies \text{¡estacionario!}
        \]
        \item Por lo que las series en diferencias ya no tienen el problema.
    \end{itemize}
\end{frame}
\begin{frame}{¿Cómo saber si $\rho = 1$? Test visual y Test de D-F}
    \begin{itemize}
        \item Naturalmente uno querría estimar el siguiente modelo para saber si la serie tiene una raíz unitaria:
        \[
        y_t = \alpha + \rho y_{t-1} + \epsilon_t
        \]
        \item Y hacer una prueba de hipótesis sobre $\rho = 1$.
        \item El problema es que si $\rho = 1$ la estimación de $\rho$ no es válida y el estadístico t de contraste no seguirá una distribución paramétrica conocida, ni siquiera en muestras grandes.
        \item En otras palabras, estaríamos usando para hacer la prueba un estadístico que no podemos estimar consistentemente y por eso la prueba no es válida.
    \end{itemize}
\end{frame}
\begin{frame}{Test de raíces unitarias}
    \begin{itemize}
        \item Dickey-Fuller (D-F) derivan distribuciones asintóticas válidas para estos casos que nos van a permitir hacer inferencia válida sobre $\rho$:
        \[
        y_t = \alpha + \rho y_{t-1} + \epsilon_t
        \]
        \[
        y_t - y_{t-1} = \alpha + \rho y_{t-1} - y_{t-1} + \epsilon_t
        \]
        \[
        \Delta y_t = \alpha + (\rho - 1) y_{t-1} + \epsilon_t
        \]
     \textcolor{red}{   \[
        \Delta y_t = \alpha + \phi y_{t-1} + \epsilon_t
        \]}
        \item Para esta especificación tenemos las tablas de distribución de D-F y por tanto podemos testear:
        \begin{itemize}
            \item $H_0: \phi = 0 \implies \rho = 1 \implies \text{Raíz unitaria (1)}$
            \item $H_1: \phi < 0 \implies \rho < 1 \implies \text{Serie estacionaria (2)}$
        \end{itemize}
        \item Construimos el estadístico t para la significación de $\phi$ y luego comparamos con los valores críticos de las tablas D-F.
    \end{itemize}
\end{frame}
\begin{frame}{Distribuciones de D-F}
    \begin{itemize}
        \item No existe una única distribución de D-F: dependerá de contra qué estemos comparando la hipótesis nula de existencia de raíz unitaria.
        \item Esto básicamente depende de cómo se comportan los datos y por tanto de cómo esté especificado el modelo.
        \item En otras palabras, en cada caso debemos pensar que nuestra hipótesis nula es que la serie tiene raíz unitaria y la alternativa puede ser:
        \begin{enumerate}
            \item Que sea estacionaria pero se mueva en el tiempo en torno a una media distinta de cero.
            \item Que sea estacionaria pero tenga una tendencia lineal en el tiempo (creciente o decreciente).
        \end{enumerate}
    \end{itemize}
\end{frame}
\begin{frame}{Distribuciones de D-F}
    \begin{itemize}
        \item D-F distinguen entre dos casos:
        \begin{enumerate}
            \item Solo constante: bajo la alternativa la serie no tiene raíz unitaria pero tiene una media distinta de cero
            \[
            y_t = \alpha + \rho y_{t-1} + \epsilon_t
            \]
            \item Constante y tendencia lineal: bajo la hipótesis alternativa la serie tiene tendencia pero es lineal en el tiempo
            \[
            y_t = \alpha + \rho y_{t-1} + \delta t + \epsilon_t
            \]
        \end{enumerate}
        \item En cada caso hay una tabla distinta con valores críticos distintos.
    \end{itemize}
\end{frame}
\begin{frame}{Valores críticos de las distribuciones de D-F}
    \begin{itemize}
        \item Los valores críticos varían según la hipótesis alternativa contra la que se quiere testear y según el tamaño de la muestra.

             \begin{figure}
        \centering
        \includegraphics[width=0.7\linewidth]{sem14/tabla.png}
    \end{figure}
        
        \item Para ver cuál es el modelo alternativo que tiene sentido, hay que empezar por mirar los datos.
    \end{itemize}
\end{frame}

\begin{frame}{Valores cr\'iticos de las distribuciones de D-F}
    Los valores críticos para la prueba de Dickey-Fuller suelen obtenerse a través de simulaciones Monte Carlo, reflejando mejor la distribución real del estadístico bajo la hipótesis nula. Esto implica que los valores críticos se derivan empíricamente y son específicos a la estructura del modelo de la prueba.
\end{frame}

%\begin{frame}{Exploración gráfica - AR(1)}
%    \begin{itemize}
%        \item AR(1) con $\rho = 0.5$
%    \end{itemize}
%             \begin{figure}
%        \centering
%        \includegraphics[width=0.7\linewidth]{sem14/images/ar1.png}
%    \end{figure}
%\end{frame}
\begin{frame}{Exploración gráfica - AR(1)}
    \begin{itemize}
        \item AR(1) con $\rho = 0.95$
    \end{itemize}
             \begin{figure}
        \centering
        \includegraphics[width=0.5\linewidth]{sem14/images/ar2.png}
    \end{figure}
    \begin{itemize}
        \item Vemos cómo aumenta la persistencia.
        \item Aquí ya podríamos pensar en diferenciar el modelo.
    \end{itemize}
\end{frame}
\begin{frame}
\frametitle{Exploraci\'on gr\'afica - I(1)}

\begin{itemize}
    \item Ar(1) con $\rho = 1$
             \begin{figure}
        \centering
        \includegraphics[width=0.7\linewidth]{sem14/images/i1.png}
    \end{figure}
    \item La persistencia es muy alta
\end{itemize}

\end{frame}
\begin{frame}
\frametitle{Exploraci\'on gr\'afica - Primera diferencia de una I(1)}

\begin{itemize}
    \item La serie con orden de integraci\'on 1 la diferenciamos, y ¿qu\'e obtenemos?
             \begin{figure}
        \centering
        \includegraphics[width=0.5\linewidth]{sem14/images/i2.png}
    \end{figure}
    \item Si la serie est\'a en logaritmos, ¿A qu\'e se parece la primera diferencia?
    \item Muchas series macroecon\'omicas siguen este tipo de procesos
\end{itemize}

\end{frame}
\begin{frame}
\frametitle{El PIB Chileno}

\begin{itemize}
    \item Queremos saber cu\'al es el orden de integraci\'on del PIB chileno.
    \item Comenzamos por una exploraci\'on gr\'afica

                 \begin{figure}
        \centering
        \includegraphics[width=0.6\linewidth]{sem14/images/azul.png}
    \end{figure}
    
   
\end{itemize}

\end{frame}
\begin{frame}
\frametitle{El PIB Chileno - Primera Diferencia}

\begin{itemize}
    \item Primera diferencia del PIB

                     \begin{figure}
        \centering
        \includegraphics[width=0.6\linewidth]{sem14/images/rojo.png}
    \end{figure}
    
    \item ¿Qu\'e impresi\'on nos da?
\end{itemize}

\end{frame}
\begin{frame}
\frametitle{El PIB Chileno - Ambas Gráficas}

\begin{itemize}
    \item Poniendo ambas gr\'aficas juntas
\end{itemize}

                 \begin{figure}
        \centering
        \includegraphics[width=0.6\linewidth]{sem14/images/azulroj.png}
    \end{figure}

\end{frame}
\begin{frame}
\frametitle{El PIB Chileno - Test de DF}

\begin{itemize}
    \item Realizamos el test de D-F para ver si se trata de una serie integrada de orden 1
    \item Dada nuestra exploraci\'on gr\'afica, corresponder\'ia usar como modelo alternativo el que incluye una tendencia temporal
$$\Delta PIB_t = \alpha + \phi PIB_{t-1} + \delta t +\epsilon_t$$
                     
    \begin{figure}
        \centering
        \includegraphics[width=0.7\linewidth]{pib.png}
    \end{figure}
    
    \item ¿Con qu\'e valor cr\'itico comparamos el estad\'istico t?
\end{itemize}

\end{frame}
\begin{frame}
\frametitle{El PIB Chileno - Test de DF}

\begin{itemize}
    \item Vamos a usar los valores de tabla de DF
    \item $\hat{\phi} = -0.0373029$ y $\hat{\sigma_{\phi}}=0.0228403$ 
    \item Como nuestra muestra es de 60, podemos usar los valores cr\'iticos de la de 50:
    \begin{itemize}
        \item VC al 5\%: -3.50
        \item $| -1.63 | < | -3.50 |$
        \item No rechazo $H_0$
    \end{itemize}
    \item Evidencia a favor de que el PIB chileno tiene una ra\'iz unitaria: no es una serie estacionaria
\end{itemize}

\end{frame}
\begin{frame}[shrink=10]
\frametitle{Conclusiones y Repaso}

\begin{itemize}
    \item Cuando trabajamos con series de tiempo debemos comenzar por ver si las series tienen tendencias temporales: determin\'isticas o estoc\'asticas
    \item Si es determinista, agregamos una tendencia:
    \[
    Y_t = \beta_0 + \beta_1 X_t + \alpha t + u_t
    \]
    \item Interpretaci\'on de $\beta_1$ ahora?
    \item Para saber si tiene una tendencia estoc\'astica (ra\'iz unitaria) procedemos a realizar el test de Dickey-Fuller
    \item Este test tiene valores cr\'iticos espec\'ificos y los mismos dependen del modelo alternativo que nos estemos planteando (con o sin tendencia determin\'istica)
    \item Para determinar el modelo alternativo comenzamos por una exploraci\'on gr\'afica
    \item Si encontramos que la serie tiene una ra\'iz unitaria (no rechazamos $H_0$) vamos a tener que diferenciar la serie y estimar nuestro modelo en diferencias.
\end{itemize}

\end{frame}


\end{document}
